\chapter{State of the art}
\label{cha:sota}

This chapter reviews the work on which the laboratory builds. It is organised by theme: AI-generated phishing and its detection (Section~\ref{sec:sota-phish}), deepfakes (Section~\ref{sec:sota-df}), adversarial examples (Section~\ref{sec:sota-adv}), the security of applications that embed LLMs (Section~\ref{sec:sota-llmsec}) and security education (Section~\ref{sec:sota-edu}). Section~\ref{sec:gap} compares existing tools and states the gap that this work addresses.

\section{AI-generated phishing and its detection}
\label{sec:sota-phish}

\subsection{Phishing as an attack vector}
Phishing remains one of the most common initial-access techniques in the threat landscape described by the European Union Agency for Cybersecurity~\cite{enisa2024threat}, and the Anti-Phishing Working Group continues to record high volumes of phishing attacks~\cite{apwg2024}. Conventional lures rely on recognisable cues: artificial urgency, generic greetings, grammatical errors, look-alike domains and requests for credentials or payment.

\subsection{Phishing generated by language models}
LLMs remove several of those cues. Hazell showed that commercially available LLMs can generate personalised spear-phishing messages for hundreds of targets at negligible cost~\cite{hazell2023spear}. Heiding et al.\ compared LLM-generated phishing emails with emails crafted by human experts in a user study, and also evaluated LLMs as detectors of malicious intent in emails~\cite{heiding2024devising}. The common conclusion is that fluent, well-tailored wording can no longer be taken as evidence that a message is legitimate.

\subsection{Machine-learning detection}
Automatic phishing detection predates LLMs. Fette et al.\ classified emails with a random forest over structural features such as IP-based URLs, link text that does not match its target, domain age and the output of a spam filter, and reported detection of more than 96\% of phishing emails with a false-positive rate of about 0.13\% on their dataset~\cite{fette2007learning}. A later systematic review by Salloum et al.\ surveys natural-language-processing approaches, ranging from TF-IDF features with classical classifiers to deep neural models~\cite{salloum2022phishing}. A TF-IDF representation with logistic regression is a standard, interpretable baseline in that body of work; it is the baseline adopted in Chapter~\ref{cha:eval}.

LLMs have also been proposed as detectors. Koide et al.\ converted emails into prompts and reported 99.70\% accuracy with GPT-4 on a dataset of 2{,}010 emails, ahead of other LLMs and conventional baselines~\cite{koide2024chatspamdetector}. Such results depend strongly on the dataset; they show that an LLM can produce a verdict and a rationale, which is the design of the email-analysis module in this work, but they do not transfer to other corpora without re-evaluation.

\subsection{Sender authentication}
If wording is no longer a reliable signal, the authenticity of the sender becomes central. Three IETF standards address it: Sender Policy Framework (SPF) authorises sending hosts for a domain~\cite{rfc7208}; DomainKeys Identified Mail (DKIM) signs messages with a domain key~\cite{rfc6376}; and Domain-based Message Authentication, Reporting, and Conformance (DMARC) lets a domain owner publish how receivers should treat mail that fails these checks~\cite{rfc9989}. These mechanisms verify the sending domain, not the content, and are therefore unaffected by how convincing an LLM makes the text.

\section{Deepfakes}
\label{sec:sota-df}

Mirsky and Lee survey the creation and detection of deepfakes of the human face and body, including reenactment and replacement techniques~\cite{mirsky2021deepfake}. Generative adversarial networks such as StyleGAN produce high-resolution synthetic faces~\cite{karras2019stylegan}. Detection research relies on benchmark datasets: FaceForensics++ provides manipulated videos produced with several methods~\cite{rossler2019faceforensics}, Celeb-DF provides higher-quality forgeries that are harder to detect~\cite{li2020celebdf}, and the DeepFake Detection Challenge released a large dataset with a public competition~\cite{dolhansky2020dfdc}.

Early detectors exploited specific artefacts. Li, Chang and Lyu observed that early face-swap videos showed abnormal eye blinking, because training images rarely contain closed eyes, and detected fakes from blink patterns~\cite{li2018ictu}. Such cues age quickly as generators improve. For synthetic speech, Müller et al.\ re-evaluated audio deepfake detectors under uniform conditions and found that their performance degrades sharply on real-world recordings outside the ASVspoof benchmark~\cite{muller2022audio}. This limited generalisation is one reason why this work does not ship a detector (Section~\ref{sec:impl-df}).

A complementary approach is provenance: the Coalition for Content Provenance and Authenticity specifies cryptographically signed metadata, known as Content Credentials, that records how a media file was created and edited~\cite{c2pa2spec}. Provenance does not detect manipulation; it allows authentic media to be verified.

\section{Adversarial examples}
\label{sec:sota-adv}

\subsection{Attacks}
Szegedy et al.\ showed that deep networks misclassify inputs that differ from correctly classified ones by small, carefully chosen perturbations~\cite{szegedy2014intriguing}. Goodfellow, Shlens and Szegedy proposed the Fast Gradient Sign Method (FGSM), a single step of size $\varepsilon$ in the direction of the sign of the input gradient of the loss~\cite{goodfellow2015explaining}:
\begin{equation}
x_{\mathrm{adv}} = \mathrm{clip}_{[0,1]}\bigl(x + \varepsilon \cdot \mathrm{sign}(\nabla_x J(\theta, x, y))\bigr).
\label{eq:fgsm}
\end{equation}
Madry et al.\ formulated robustness as a min--max problem and used Projected Gradient Descent (PGD), an iterated FGSM with projection back onto the $\varepsilon$-ball, as a strong first-order attack~\cite{madry2018towards}. Biggio and Roli review a decade of adversarial machine learning, including evasion and poisoning attacks~\cite{biggio2018wild}. The Torchattacks library implements these attacks for PyTorch~\cite{kim2020torchattacks}.

\subsection{Evaluation and defences}
Evaluating robustness is error-prone. Carlini et al.\ set out principles for sound evaluations, including a precise threat model, strong adaptive attacks and sanity checks such as accuracy approaching zero for large budgets~\cite{carlini2019evaluating}. Input transformations such as JPEG compression, bit-depth reduction and total-variance minimisation have been proposed as defences~\cite{guo2018countering}; however, Athalye, Carlini and Wagner showed that defences relying on obfuscated gradients can be circumvented by adapted attacks, defeating six of seven such defences completely~\cite{athalye2018obfuscated}. The adversarial module of this work therefore presents preprocessing only as a defence to be tested, not as a guarantee.

\subsection{Model choice}
Residual networks~\cite{he2016resnet} are a standard architecture in this literature. CIFAR-10 consists of $32 \times 32$ colour images in ten classes~\cite{krizhevsky2009cifar}. Pretrained CIFAR-10 residual networks are distributed in the \texttt{pytorch-cifar-models} repository~\cite{chenyaofo2021cifar}, which provides the model used in this work.

\section{Security of LLM-integrated applications}
\label{sec:sota-llmsec}

An application that inserts user text into an LLM prompt inherits the risk of prompt injection. Perez and Ribeiro showed that simple handcrafted inputs can hijack a model's goal or leak its prompt~\cite{perez2022ignore}. Greshake et al.\ introduced indirect prompt injection, in which instructions are hidden in data that the application later passes to the model rather than in the user's own request, and demonstrated it against deployed systems~\cite{greshake2023signed}. An email-analysis service is an instance of the indirect case, because the email body is chosen by a third party. These works motivate the input handling described in Section~\ref{sec:limits}.

\section{Security education}
\label{sec:sota-edu}

Research on anti-phishing education indicates that interactive training can outperform passive material. Sheng et al.\ evaluated Anti-Phishing Phil, an online game, in a 42-participant laboratory study and found that participants who played the game identified phishing websites better than those who read existing online training material~\cite{sheng2007phil}. Kumaraguru et al.\ studied anti-phishing training designs and their evaluation in more depth~\cite{kumaraguru2010johnny}. Practitioner platforms such as TryHackMe and Hack The Box offer hands-on security exercises across many topics~\cite{tryhackme,hackthebox}. These platforms are broad and hosted by third parties; they are not small applications that a student can deploy, inspect and explain in full.

\section{Comparison and gap}
\label{sec:gap}

Table~\ref{tab:gap} compares representative tools and resources against three criteria derived from the problem statement in Section~\ref{sec:problem}: coverage of the three risks, availability as a self-hostable teaching application, and transparency about whether an output is measured, generated or configured.

\begin{table}[htbp]
\centering
\caption{Comparison of related tools and resources.}
\label{tab:gap}
\scriptsize
\begin{tabular}{p{2.6cm}p{1.7cm}p{1.7cm}p{1.7cm}p{1.5cm}p{1.8cm}}
\toprule
Tool or resource & Phishing & Deepfakes & Adversarial & Self-hostable teaching app & Output provenance shown \\
\midrule
GoPhish~\cite{gophish} & Campaigns & --- & --- & Yes & Not applicable \\
FaceForensics++, Celeb-DF~\cite{rossler2019faceforensics,li2020celebdf} & --- & Datasets & --- & No & Ground-truth labels \\
Torchattacks~\cite{kim2020torchattacks} & --- & --- & Library & No (library) & Computed \\
Anti-Phishing Phil~\cite{sheng2007phil} & URL training & --- & --- & Game & Not applicable \\
TryHackMe, Hack The Box~\cite{tryhackme,hackthebox} & Some rooms & Rare & Rare & No (hosted) & Varies \\
ChatSpamDetector~\cite{koide2024chatspamdetector} & LLM detection & --- & --- & Research tool & LLM rationale \\
\textbf{InsightFlow AI (this work)} & Generate and analyse & Scenarios (no media) & FGSM/PGD & Yes & Per response \\
\bottomrule
\end{tabular}
\end{table}

Each risk is well covered on its own: phishing training and detection, deepfake datasets and detectors, and adversarial-attack libraries are all mature. No resource identified in the search combines the three risks in one self-hostable application for teaching, and the general-purpose platforms do not distinguish, per response, between a computed result, an LLM's narration and a value configured by the tool's author. The gap addressed by this work is therefore one of integration and transparency for teaching, not a new attack or detection technique.

\section{Summary}
\label{sec:sota-summary}

LLMs make phishing wording an unreliable signal, which shifts attention to sender authentication; deepfake detectors generalise poorly, which favours provenance and awareness; and adversarial robustness must be evaluated with strong attacks and verified budgets. Chapter~\ref{cha:req} turns these observations into requirements for the laboratory.
